\documentclass{article}

\usepackage{amsmath,amssymb}
\usepackage{xurl}
\usepackage[hidelinks]{hyperref}
\setlength{\emergencystretch}{2em}

% Shared commands for notes in docs/paper-gaps/.

\DeclareUnicodeCharacter{2080}{\ifmmode{}_{0}\else\textsubscript{0}\fi}
\DeclareUnicodeCharacter{2081}{\ifmmode{}_{1}\else\textsubscript{1}\fi}
\DeclareUnicodeCharacter{2082}{\ifmmode{}_{2}\else\textsubscript{2}\fi}
\DeclareUnicodeCharacter{2099}{\ifmmode{}_{n}\else\textsubscript{n}\fi}

\newcommand{\Lean}{\textsc{Lean}}
\ExplSyntaxOn
\NewDocumentCommand{\leanid}{m}{
  \ifmmode
    \textsf{\detokenize{#1}}
  \else
    \group_begin:
    \urlstyle{sf}
    \tl_set:Nn \l_tmpa_tl {#1}
    \tl_replace_all:Nnn \l_tmpa_tl {\_} {_}
    \exp_args:NV \path \l_tmpa_tl
    \group_end:
  \fi
}
\ExplSyntaxOff
\providecommand{\tr}{\operatorname{tr}}
\newcommand{\tnleanIssueBase}{https://github.com/LionSR/TNLean/issues/}
\newcommand{\tnleanPrBase}{https://github.com/LionSR/TNLean/pull/}
\newcommand{\tnleanDocBase}{https://sirui-lu.com/TNLean/docs/}
\newcommand{\ghissue}[1]{\href{\tnleanIssueBase#1}{GitHub issue \##1}}
\newcommand{\ghpr}[1]{\href{\tnleanPrBase#1}{PR \##1}}
\newcommand{\doclink}[1]{\href{\tnleanDocBase#1.html}{\path{#1}}}

% Dirac notation and the identity operator, as in blueprint/src/macros/common.tex.
% \providecommand keeps note-local definitions authoritative.
\usepackage{dsfont}
\providecommand{\ket}[1]{|#1\rangle}
\providecommand{\bra}[1]{\langle #1|}
\providecommand{\braket}[2]{\langle #1 | #2 \rangle}
\providecommand{\ketbra}[2]{|#1\rangle\!\langle #2|}
\providecommand{\Id}{\mathds{1}}
\providecommand{\one}{\mathds{1}}
\providecommand{\C}{\mathbb{C}}
\providecommand{\MN}[1]{M_{#1}(\C)}
\providecommand{\MD}{M_D(\C)}

% Verdict marker: \gapnote{<kind>}{<status>}. Kind is the policy
% classification (clarification, local-correction, scope-restriction,
% unfaithful, false-source, open-gap); status is open, wip (elimination
% underway), resolved, or historical. Machine-read by the site index;
% prints a verdict line.
\newcommand{\gapnote}[2]{\par\noindent\textbf{Verdict:} #1 (#2).\par}


% Fallbacks keep this note compilable when a different command.tex is found first.
\providecommand{\leanid}[1]{\textsf{\detokenize{#1}}}
\providecommand{\gapnote}[2]{\par\noindent\textbf{Verdict:} #1 (#2).\par}
\providecommand{\ket}[1]{|#1\rangle}
\providecommand{\bra}[1]{\langle #1|}

\title{Bounded-Change Expansions of Births, Deaths and Exchanges}
\author{}
\date{2026-10-08}

\begin{document}
\maketitle
\gapnote{scope-restriction}{resolved}

\section{The source lemmas}

Lemmas~6.5 and~6.6 of the polynomial PEPS approximation manuscript of
September~24, 2026 (\path{05-frames.tex}, lines~396--570) change the ownership
of registers of encoded frames.

\emph{Lemma~6.5 (homogeneous birth and death).} In an encoded frame
$\mathcal F$, let $T$ be a set of raw sites owned by a party $P_\circ$ and
disjoint from every outer hole footprint, and put
$E=\{x:\text{the raw owner of }x\neq P_\circ\}\cup H^+(\mathcal F)$ and
$U=\Lambda\setminus(T\cup E)$. If $I_\Omega(T:E)\le L^{-60}$, changing the owner
of $T$ to $Q_\circ$, with every hole, encoding and tag owner kept, ``is a bounded
change involving only $P_\circ,Q_\circ$, with reference-vector error at most
$L^{-30}$''; the reversed death obeys the same bound when the conditions hold in
the frame after the death.

\emph{Lemma~6.6 (two-sheet exchange).} For two encoded frames and a region $Y$
such that every hole's outer square lies on one side of $\partial Y$, with
$Z=\{x:\text{the two old raw owners of }x\text{ are not both }P_\circ\}\cup
H^+(\mathcal F_1)\cup H^+(\mathcal F_2)$, $T=Z\cap Y$, $E=Z\setminus Y$ and
$U=\Lambda\setminus Z$: if $I_\Omega(T:E)\le L^{-60}$, exchanging the raw
ownership assignments and the encoded holes inside $Y$ ``can be implemented by
private contractions and register renaming, with reference-vector error at most
$4L^{-30}$'', using $P_\circ$ and the owners of the renamed registers and tags,
hence is a bounded change in the sense of Theorem~5.2.

In both lemmas the information bound is a hypothesis; it is supplied later by the
manuscript's proposition \texttt{prop:angular}.

\section{What the formal statements cover}

Both lemmas are formalized with exactly their source hypotheses, the information
bound included.\footnote{%
\leanid{TNLean.PEPS.EncodedFrame.Frame.birth},
\leanid{TNLean.PEPS.EncodedFrame.Frame.death},
\leanid{TNLean.PEPS.EncodedFrame.TwoSheetExchange.exchange} and
\leanid{TNLean.PEPS.EncodedFrame.TwoSheetExchange.exchange_ofFrames}.}
The formal statements produce the canonical maps of the proofs and prove:
\begin{enumerate}
\item for the birth and the death, that the canonical map
  $B=\mathcal V^\dagger(\ket s\bra s_{T\mathcal B_T}\otimes\mathrm{id}_{E\mathcal
  B_E})\mathcal V$ is a Hermitian contraction, that it acts as the identity on
  $E$ and on all tags, that it commutes with every hole encoder, and that the
  reference error is at most $L^{-30}$; the map factors as
  $B=(\mathcal V^\dagger S)(S^\dagger\mathcal V)$, where $S$ inserts the single
  normalized vector $s$ on $T\mathcal B_T$, and every register on which $B$ acts
  lies outside $E$, so it is held by $P_\circ$ before the birth and by $P_\circ$
  or $Q_\circ$ after it; since the frames before and after differ only in raw
  owners, their encodings and reference vectors coincide, so the change of owner
  itself is carried entirely by the bounded-change clause below;
\item for the exchange, the exact identity $\mathcal RK_{\rm in}=K_{\rm out}F_{Y\cap
  \Lambda}$ for the renaming $\mathcal R$, the identity $D_UF_T=(\mathcal
  V^{\otimes2})^\dagger F_{T\mathcal B_T}\mathcal V^{\otimes2}$, the identity
  $CK_{\rm in}=K_{\rm out}D_UF_T$ for $C=D_UF_A\mathcal R$, contractivity of $C$,
  and the error bound $4L^{-30}$ on the product of the two encoded reference
  vectors, which are not assumed normalized; the corrections $D_U$ and $F_A$
  act on the registers of $U$ alone, which are held by $P_\circ$ on both sheets
  before and after the exchange, and the renaming keeps every register at its
  party.
\end{enumerate}

The bounded-change clauses are formalized on the registers of the frames, as
the source has them: one raw register $\mathbb C^q$ per site and one register
per tag, each held by its owner (\path{05-frames.tex}, lines~15--20 and
Definition~6.1), with their tensor product identified with the canonical
coordinates of the tags and the sheet in which the encodings
act.\footnote{%
\leanid{TNLean.PEPS.EncodedFrame.layoutRegs},
\leanid{TNLean.PEPS.EncodedFrame.layoutIso},
\leanid{TNLean.PEPS.EncodedFrame.Frame.regs},
\leanid{TNLean.PEPS.EncodedFrame.Frame.regIso}.}
A monomial is a composition of contractions acting on registers of one party,
normalized pair sources and normalized pair effects on registers of two named
parties, and exchanges of tensor factors, as for Lemma~5.1.
\begin{enumerate}
\item The birth is an allowed monomial from the registers of $\mathcal F$ to
  those of the frame in which $T$ is owned by $Q_\circ$, using only $P_\circ$
  and $Q_\circ$ and at most one normalized pair source, whose operator in
  canonical coordinates is $\mathrm{id}_{\rm tags}\otimes B$; its
  reference-vector error, read on the registers, is at most $L^{-30}$. It
  moves the registers of $T$ and of $U$ in front by exchanges of tensor
  factors, groups each block into one register by a unitary at its owner
  $P_\circ$, applies the private contraction $\bra s(\mathrm{id}_T\otimes V)$
  at $P_\circ$, the pair source $s$ on $T$ held by $Q_\circ$ and $\mathcal
  B_T$ held by $P_\circ$, and the private contraction $V^\dagger$ at
  $P_\circ$, and ungroups $T$ at $Q_\circ$ and $U$ at $P_\circ$. For
  $Q_\circ=P_\circ$ the middle step is the single private contraction
  $K^\dagger K$.\footnote{%
\leanid{TNLean.PEPS.EncodedFrame.Frame.birth_bounded}.}
\item The death is the same with one normalized pair effect $\bra s$ on $T$
  held by $Q_\circ$ and $\mathcal B_T$ held by $P_\circ$ and no pair
  source.\footnote{%
\leanid{TNLean.PEPS.EncodedFrame.Frame.death_bounded}.}
\item The exchange is the renaming $\mathcal R$, written as a word of
  exchanges of tensor factors that splits the tags of each frame into those of
  the holes outside and inside $Y$, moves the raw registers of the sites in
  $Y$ next to the inside tags, exchanges these blocks between the sheets and
  merges back, followed by one private contraction at $P_\circ$ on the raw
  registers of $U$ of both sheets. It uses only $P_\circ$ and no pair
  resource, its operator in canonical coordinates is $C=D_UF_A\mathcal R$,
  and its reference-vector error is at most $4L^{-30}$. For frames whose holes
  are listed in any order, the canonical identification of tag orderings is
  itself a word of exchanges of tensor factors, and the composite has operator
  $CG$.\footnote{%
\leanid{TNLean.PEPS.EncodedFrame.TwoSheetExchange.exchange_bounded},
\leanid{TNLean.PEPS.EncodedFrame.TwoSheetExchange.exchange_ofFrames_bounded}.}
\end{enumerate}
The renaming moves no register to another party. The source's assertion that
the exchange uses $P_\circ$ and the owners of the renamed registers and tags is
satisfied with the stronger property that only $P_\circ$ is used, since
exchanges of tensor factors involve no party.

The formal frames record the outer footprints of holes by their physical
samples. The source's condition that every outer square lies on one side of
$\partial Y$ enters as a classification of the holes as outside or inside $Y$,
which is part of the data, together with its consequence for the samples: the
sample of a hole classified as outside avoids $Y$, and that of a hole
classified as inside lies in $Y$. A hole whose outer square has empty sample is
compatible with both sides, so the side read from the plane must be supplied,
and the formal statement accepts it. This is a representation of the source
hypothesis, not an additional one.

\section{Resolution}

An earlier version of the formal statements proved the bounded-change clauses
only on a layout in which the raw registers of $T$, and those of $U$, were
grouped into one register each, followed by an arbitrary list of untouched
registers, and identified the renaming of the exchange with an identification
of tensor factors rather than a composition of exchanges. The list of the
registers of a frame, one per site and per tag, and its identification with
the canonical coordinates are now constructed, together with the words that
move and group registers and the words that implement the renaming and the
corrections. Composing the grouped monomials with these words gives the
statements above, with the hypotheses of Lemmas~6.5 and~6.6. No restriction
remains.

\end{document}
